\section*{Sets and Indices}

There are no nontrivial index sets in the implemented model. The formulation uses aggregate variables for:
\begin{itemize}
    \item product A,
    \item product B,
    \item by-product C sales and disposal,
    \item process stage 1 and process stage 2 operating mode.
\end{itemize}

\section*{Parameters}

All parameter names below match the data field names in \texttt{general\_parameters.csv}.

\begin{itemize}
    \item $t^{A}_{1}$ (\texttt{product\_a\_process\_1\_time}): process 1 time per unit of product A.
    \item $t^{A}_{2}$ (\texttt{product\_a\_process\_2\_time}): process 2 time per unit of product A.
    \item $t^{B}_{1}$ (\texttt{product\_b\_process\_1\_time}): process 1 time per unit of product B.
    \item $t^{B}_{2}$ (\texttt{product\_b\_process\_2\_time}): process 2 time per unit of product B.
    \item $T_{1}$ (\texttt{available\_time\_process\_1}): available process 1 time in normal mode.
    \item $T_{2}$ (\texttt{available\_time\_process\_2}): available process 2 time in normal mode.
    \item $T^{H}_{1}$ (\texttt{available\_time\_process\_1\_high}): available process 1 time in high-throughput mode.
    \item $T^{H}_{2}$ (\texttt{available\_time\_process\_2\_high}): available process 2 time in high-throughput mode.
    \item $\gamma$ (\texttt{byproduct\_c\_per\_unit\_b}): units of by-product C generated per unit of product B.
    \item $\bar S^{C}$ (\texttt{max\_byproduct\_c\_sales}): maximum by-product C sales.
    \item $\tau^{C}$ (\texttt{threshold\_c}): disposal threshold for C at base disposal cost.
    \item $M^{C}$ (\texttt{bigM\_disposal}): upper bound used on total C disposal.
    \item $\pi_A$ (\texttt{profit\_per\_unit\_a}): profit per unit of product A.
    \item $\pi_B$ (\texttt{profit\_per\_unit\_b}): profit per unit of product B.
    \item $\pi_C$ (\texttt{profit\_per\_unit\_c}): profit per unit of by-product C sold.
    \item $d_L$ (\texttt{disposal\_cost\_per\_unit\_c}): base disposal cost per unit of C.
    \item $d_H$ (\texttt{disposal\_cost\_per\_unit\_c\_high}): higher disposal cost per unit of C beyond the threshold.
\end{itemize}

\section*{Decision Variables}

\begin{itemize}
    \item $x_A$ (code: \texttt{xA}) $\ge 0$: units of product A produced.
    \item $x_B$ (code: \texttt{xB}) $\ge 0$: units of product B produced.
    \item $s_C$ (code: \texttt{cSold}) $\ge 0$: units of by-product C sold.
    \item $u_C$ (code: \texttt{cDispLow}) $\ge 0$: units of by-product C disposed at base cost.
    \item $v_C$ (code: \texttt{cDispHigh}) $\ge 0$: units of by-product C disposed at higher cost.
    \item $z^{H}_{1}$ (code: \texttt{z1\_high}) $\in \{0,1\}$: 1 if process stage 1 uses high-throughput mode.
    \item $z^{H}_{2}$ (code: \texttt{z2\_high}) $\in \{0,1\}$: 1 if process stage 2 uses high-throughput mode.
\end{itemize}

\section*{Objective}

The implemented model maximizes total profit:
\[
\max \;
\pi_A x_A + \pi_B x_B + \pi_C s_C - d_L u_C - d_H v_C .
\]

\section*{Constraints}

\paragraph{Process stage capacities with mode selection.}
\[
t^{A}_{1} x_A + t^{B}_{1} x_B \le T_{1} + (T^{H}_{1} - T_{1}) z^{H}_{1},
\]
\[
t^{A}_{2} x_A + t^{B}_{2} x_B \le T_{2} + (T^{H}_{2} - T_{2}) z^{H}_{2}.
\]

\paragraph{High-throughput mutual exclusion.}
At most one reaction stage may use the high-throughput setting:
\[
z^{H}_{1} + z^{H}_{2} \le 1.
\]

\paragraph{By-product C balance.}
\[
s_C + u_C + v_C = \gamma x_B.
\]

\paragraph{Sales limit for by-product C.}
\[
s_C \le \bar S^{C}.
\]

\paragraph{Threshold-based disposal structure for by-product C.}
The amount disposed at base cost is capped by the threshold, and total disposal is capped by the supplied big-$M$ bound:
\[
u_C \le \tau^{C},
\]
\[
u_C + v_C \le M^{C}.
\]

\paragraph{Variable domains.}
\[
x_A, x_B, s_C, u_C, v_C \ge 0,
\qquad
z^{H}_{1}, z^{H}_{2} \in \{0,1\}.
\]

\section*{Notes on Implementation}

\begin{itemize}
    \item The model in \texttt{current\_heuristic.py} is a MILP solved with Gurobi through the PuLP-compatible interface, not by exact enumeration or dynamic programming.
    \item The high-throughput rule is implemented as a binary gate on each process-stage capacity: each stage gets either its normal capacity or its high-throughput capacity, and the mutual-exclusion constraint enforces that at most one stage can receive the high-throughput capacity in the period.
    \item The disposal rule for by-product C is implemented exactly as in code using two disposal variables. The first segment $u_C$ is limited by \texttt{threshold\_c}, and any additional disposal may go to $v_C$ at the higher unit cost.
    \item The constraint $u_C + v_C \le M^{C}$ uses the data field \texttt{bigM\_disposal} as an explicit upper bound on total disposal. This is part of the implemented model and should not be removed or tightened unless the code is changed.
    \item Production variables are continuous in the implementation; there are no integrality restrictions on $x_A$ or $x_B$.
\end{itemize}
